\begin{table}[htbp]\centering\small\caption{Exploratory annotation case: QASPER. Unit indices are zero-based in the released mapping; source-unit IDs use paper:paragraph:index or abstract:sentence:index. Brief excerpts are from the identified source. No new human semantic validation is claimed.}
\begin{tabular}{>{\raggedright\arraybackslash}p{79pt}>{\raggedright\arraybackslash}p{329pt}}
\toprule
{Field} & {Recorded example} \\
\midrule
{Source IDs} & {1909.12208; item 14fdc8087f2a62baea9d50c4aa3a3f8310b38d17} \\
{Question / claim} & {What supports the claim that enhancement in training is advisable as long as enhancement in test is at least as strong as in training?} \\
{Answer / label} & {Paraphrase: an extensive experimental evaluation on difficult CHiME-5 dinner-party data.; Paraphrase: training enhancement improves accuracy provided it is no stronger than test enhancement.} \\
{Evidence family} & {$E_{1}$=\{17,22,23,34\}; $E_{2}$=\{23\}} \\
{Exact prefix} & {[34,1,2,23,3] at k=5} \\
{Completion} & {H=1, C=1, A=0; $d_{C}$=4, $d_{A}$=23} \\
{Brief unit excerpts} & {17: ``Experiments were performed using the CHiME-5 data.''; 22: ``An extensive set of experiments''; 23: ``the recognition accuracy improves monotonically''; 34: ``CHiME-5 dinner party data''} \\
{Interpretation} & {The singleton \{23\} is nested within \{17,22,23,34\}; completing the singleton makes C=1 before the full union is covered. The two recorded normalized answers differ. This illustrates annotation nesting and answer-text variation, not same-answer alternative sufficient evidence.} \\
\bottomrule\end{tabular}\end{table}
